% !TeX root = ../papers/nc-p3.tex
\section{正の固定曲線の通常階数15窓の全点分類}
\label{sec:rank-fifteen-all-ordinary}
\begin{theorem}\label{thm:rank-fifteen-all-ordinary}
標数零で定理\ref{thm:rank-fifteen-covers}の正の固定曲線から来る二表示の
全ての通常階数15窓は、標準14641型quadratic Koszul左右AS正則である。
四始点で共通の四つの中心四次式を選び、各始点の経路をその多項式環上の
自由階数256の加群にできる。従ってこの枝に通常な非正則点は存在しない。
定理\ref{thm:rank-fifteen-formal-boundary}の四点を加えた規格化窓の添字保持同型類は
\[
        y\ne0,\qquad p^2=x^3-27y,\qquad q^2=(1-x)^2-4y
\]
の各点に一対一に対応する。一位置移動は $(p,q)\mapsto(-p,-q)$ で、
$x=3,3/4$ で $p=q=0$ の二点だけが周期1、その他は全て周期2である。
\end{theorem}
\begin{proof}[九次以後の四つの上界]
第\ref{sec:three-center-pencil}節までの帰着と
定理\ref{thm:rank-fifteen-ninth-character}により、剰余標数3の
$X\notin\mathbb F_9$, $Y\ne0,(1-X)^2$ の単元対で、長さ10から13の
始点0の上界 $10,4,1,0$ を示せばよい。$c^2=X$ の符号を指定し、
$B=Y/X^3$、$(v^3-v)^2=B$ の六表示を使う。

命題\ref{prop:ordinary-bivariate-border-cut}と同じく第2,3,4,5列を課し、
第1列を変数 $Z$ に送る。初期の五中心商の主元は通常通り記録する。
長さ9以後では、右倍に用いるのは独立行を選ぶ前の元の行である。
よって後の次数の上界を証明するとき、それ以前に独立と判定した残余の
独立性を特殊化先で仮定する必要はない。前段の選択が従属に変わっても、
選ばれた各行とその右倍は依然として真の関係である。

現在の文字の段でも、選んだ原行と主元列でできる一つの小行列式だけを残す。
その値は、行と列の順序の符号を除き、有理係数としての主元の積である。
実際、選ばれた行だけを同じ順で消去する行加算はその小行列式を変えない。
原始多項式行に付随する有理倍率を含めて主元を掛け、約分後の分子を用いる。
特殊化でこの分子が単元なら、初期表が定義される限り、途中の個々の消去を再現できなくても
必要な原行の独立性が保証される。十一、十三次はWalsh文字の平行移動と
定理\ref{thm:rank-fifteen-ninth-character}の四平面の開条件により、
文字0だけでよい。偶数次は四文字の小行列式を同じ表示で全て使う。

各次数に必要な因子だけを $(v,c)$ へ引き戻した平方自由積を $D$ とする。
$D$ を $1,v,\ldots,v^5$ に展開した六係数の共通因子は、四行とも次の同じ多項式 $G(c,B)$
である。$u=cv$ として
\[
  F(c,v)=u^6-(c^2+1)u^4+(c^4-c^3+c)u^2+c^3-c^2-c+1,\qquad
  G(c,(v^3-v)^2)=\prod_{j\in\mathbb F_3}F(c,v+j).
\]
右辺の六係数は定数係数以外が零で、$G$ の二次数は $(18,3)$ である。
$G$ を除いた六係数に符号を固定した終結式上界を適用すると次の表を得る。
\[
\begin{array}{c|c|c|c|c}
 \text{長さ}&\text{検査する文字}&(\deg_v D,\deg_c D)&
  c\text{ の終結式上界の次数}&\text{残余軌道数}\\ \hline
 10&0,1,2,3&(120,114)&6&2\\
 11&0&(84,80)&0&0\\
 12&0,1,2,3&(171,170)&9&3\\
 13&0&(88,86)&3&1
\end{array}
\]
全四行の例外を合わせると三軌道だけである。
$c^3+2c^2+c+1=0$ で $B=1+c+2c^2$、
$c^3+2c^2+2c+2=0$ で $B=1$ または $2+2c+c^2$ となる。
各軌道は元の四窓に戻して全四始点の長さ13までを直接検査する。
\path{computations/verify_rank_fifteen_all_ordinary.py} は、二変数の語関係から
四行の小行列式・六係数・共通曲線・三軌道の検査を再現する。
\end{proof}

\begin{proof}[共通例外曲線の全点]
$G(c,B)=0$ なら、所定の $c$ を変えず、六つの $v$ のうち
$F(c,v)=0$ を満たすものが取れる。これは上の積の恒等式に従う。
$F$ は $\mathbb F_3[c,v]$ で既約であり、$c\ne0$ で $u=cv$ を使うと表示した
モニック六次式になる。$\mathbb F_3(c)[u]/(F)$ 上でまず $v=u/c$ の表示で第1列を落とし、
次に $v'=u/c+1$ の表示で第4列を落とす。いずれも始点0の四中心商は長さ13の消滅まで標準次元を持つ。

除外点の処理では、除算のノルムの平方自由因子を $c$ の一変数多項式として記録する。
具体的には $t=u^2$ として元を $(a(t)+u b(t))/d(c)$、
$\deg_t a,\deg_t b\le2$ と書く。逆元は $a^2-t b^2$ の三次拡大での乗法行列の余因子と
行列式で計算できる。二表示の記録の次数は6986と17641である。その最大公約元から
$c^{18}-c^2$ の根を除くと、残りは平方自由な次数237で、既約因子の数は
次数 $4,5,6,7,8$ の順に $5,11,15,8,2$ である。
各因子の一根 $c$ で $F(c,v)$ を有限体上分解し、全ての $B=(v^3-v)^2$ を
$\mathbb F_3(c)$ のFrobenius軌道に分ける。
$B\in\mathbb F_9$ は定理\ref{thm:rank-fifteen-cubic-pencils}、
$Y/(1-X)^2\in\mathbb F_9$ は既出の放物線束の定理で覆われる。
残る100軌道では、所定の $c$ の六表示と規格化五列の四中心から選び、
最初の二次関係に戻って十三次までの上界を検算する。これらの計算は
\path{computations/verify_rank_fifteen_signed_curve.py} にあり、
除算する六次拡大の算術は \path{computations/rank_fifteen_signed_curve_field.py} にある。

これで、各次数で、共通曲線の上も外も含めて、始点0で所定の上界を持つ表示が得られた。
必要な開集合の交叉は奇数次の四文字については一つの剰余表示の中で行う。
五つの中心を持ち上げ、指定した符号 $\bar p=c^3$ で目的の標数零の窓へ運ぶ。
四つの次数で同じ表示を選ぶ必要はない。各移動後の窓にも適用すると、元の窓の
全四次中心のGrassmann多様体上で各始点・各次数について非空開集合が得られる。
九次までに得た開集合も合わせて交差させ、同一の四中心で長さ13まで全ての上界が成立する。
命題\ref{prop:klein-central-strata}の全次数の自由性と左右Koszul完全性、有限双対の
完全対合から所望の正則性を得る。
最後に定理\ref{thm:rank-fifteen-global-invariants}と
\ref{thm:rank-fifteen-formal-boundary}の同型条件・通常表示の像と境界四点を適用する。
同型不変量の曲面は
$(2p+3\sqrt3q)(2p-3\sqrt3q)=(x-3)^2(4x-3)$
と書ける既約有理曲面の開部分で、一位置移動の固定点は述べた二点だけである。
\end{proof}
